\documentclass[10pt,a4paper]{amsart}
\usepackage[margin=19mm]{geometry}
\usepackage{amsmath,amssymb,mathtools}
\usepackage[scr=rsfs,cal=euler]{mathalfa}
\usepackage{xfrac}
\usepackage[unicode,hidelinks,bookmarks=false]{hyperref}
\newcommand{\ie}{i.e.\ }
\newcommand{\cf}{cf.\ }
\newcommand{\resp}{respectively\ }
\newcommand{\Subsection}{\subsection}
\providecommand{\nobreakdash}{\nobreak\hbox{-}}
\setlength{\emergencystretch}{2em}
\multlinegap0pt
\usepackage{amssymb}
\usepackage{xcolor}
\usepackage{tikz}
\newtheorem{theorem}{Theorem}
\newtheorem{corollary}{Corollary}
\newtheorem{proposition}{Proposition}
\title{Schur polynomials twisted by roots of unity and reciprocal pairs: exactly three factors, and where they vanish}
\author{Carles Mar\'in}
\date{Source: arXiv:2608.09619v3 (CC BY 4.0)}
\begin{document}
\maketitle

\noindent\emph{Source and licence.} Schur polynomials twisted by roots of unity and reciprocal pairs: exactly three factors, and where they vanish. Version arXiv:2608.09619v3 (CC BY 4.0), distributed by arXiv under the Creative Commons Attribution 4.0 International licence, http://creativecommons.org/licenses/by/4.0/, as recorded in the arXiv licence metadata for this exact version and shown on the abstract page https://arxiv.org/abs/2608.09619v3. Full licence notice: \texttt{licenses/arxiv-2608.09619-CC-BY-4.0.txt}.\par\smallskip

\noindent\textit{Editorial scope.} Counted unit: the article's proposition prop:A1gen (source0.tex lines 3253--3256). The article's complete original proof of it is delivered with it (source0.tex lines 3258--3292, one \texttt{proof} environment of 35 lines). No mathematics is written, repaired or re-derived: the statement and the proof are the article's own lines, in the article's own notation, and no retained line is substituted or re-flowed.  Retained uncounted as the source-local context the proof needs: lines 2903--2913; lines 3191--3198.  Nothing was omitted from any retained span, so the package carries no omission record.  No retained line contains a citation, so the wrapper carries no bibliography and no bibliography entry.  No retained line contains a figure environment, an image-inclusion command or drawing code, so no image is delivered and the package declares no asset dependency.  Editorial formatting only, in addition: an AMS \texttt{amsart} preamble that restates the article's own declarations and macros used by the retained lines, namely the environments \texttt{theorem}, \texttt{corollary}, \texttt{proposition} on separate counters, together with the standard packages those lines need.  The retained environments print in this document as Theorem 1, Corollary 1, Proposition 1 in order of appearance, whereas the article prints the same items with the numbering of its own full document.  The complete native source of the article as distributed in the arXiv e-print for this exact version is bundled unchanged as \texttt{sources/207258/source0.tex}.
\begin{theorem}[the extremal set]\label{thm:Ggen}
Assume $n_i\ge1$ for every $i$. Then
\begin{enumerate}
\item[\rm(a)] $g\in G$ if and only if the multiset of increments it selects, in the sense of Step 2
below, consists of all increments $>\tau$ together with enough increments equal to $\tau$;
\item[\rm(b)] $|G|\le2$, and $|G|=2$ if and only if $\Delta_{(r)}=\Delta_{(r+1)}$, in which case $t$
is even, the two tied increments belong to classes $i$ and $i+t/2$ --- both necessarily in
$\mathcal{E}$ --- and the two maximisers differ exactly in those two classes;
\item[\rm(c)] if $t$ is odd then $|G|=1$.
\end{enumerate}
\end{theorem}

\begin{corollary}[the reflection]\label{cor:reflectgen}
If $\bigl[\det(x_i^{\beta_j})\bigr]_{D_1}=0$ then $|G|=2$ and $T_{\mathrm b}=\tau-T_{\mathrm a}$.
Consequently $\tau-\mathcal{K}=\mathcal{K}$ and, since $x\mapsto\tau-x$ also fixes $\{p_1,p_2\}$ and
$\{q_1,q_2\}$ setwise,
\begin{equation}\label{eq:Ssymgen}
  \tau-\bigl(\mathcal{K}\cup\{p_1,p_2,q_1,q_2\}\bigr)=\mathcal{K}\cup\{p_1,p_2,q_1,q_2\}.
\end{equation}
\end{corollary}

\begin{proposition}[the extremes avoid the common part]\label{prop:A1gen}
Suppose $\bigl[\det(x_i^{\beta_j})\bigr]_{D_1}=0$. Then
$\max\mathcal{S}\notin g_{\mathrm{com}}$ and $\min\mathcal{S}\notin g_{\mathrm{com}}$.
\end{proposition}

\begin{proof}
By Corollary \ref{cor:reflectgen}, $T_{\mathrm b}=\tau-T_{\mathrm a}$. Since $x\mapsto\tau-x$ reverses
order, it carries the bottom half of $T_{\mathrm a}$ to the top half of $T_{\mathrm b}$ and back:
\begin{equation}\label{eq:halvesgen}
  H_{\mathrm b}=\tau-L_{\mathrm a},\qquad L_{\mathrm b}=\tau-H_{\mathrm a}.
\end{equation}
Recall from Step 2 of the proof of Theorem \ref{thm:Ggen} that a maximiser is described by the
integers $k_i$: inside class $i$, read decreasingly, the first $k_i$ excess values lie in the top half,
the next one is the omitted $g_i$, and the remaining ones lie in the bottom half. Recall also from
Theorem \ref{thm:Ggen}(a) that a maximiser selects every increment $>\tau$, that the increments of a
class are selected in prefix order, and that the only increments equal to $\tau$ are the two tied ones,
which lie in the two tied classes.

\emph{The maximum.} Put $m=\max\mathcal{S}$ and let $i$ be its class, so $m=c_{i,1}$ and $n_i\ge2$;
write $c=c_{i,2}$. Then $m\in g_{\mathrm a}\cap g_{\mathrm b}$ if and only if $k_i=0$ in both
maximisers. Assume it. Then class $i$ is not tied: a maximiser taking the tied increment of a tied
class has $k_i\ge1$ there. Hence no increment of class $i$ is $\ge\tau$, and in particular
\[
  \Delta_i(1)=m+c<\tau .
\]
On the other hand $k_i=0$ places $c$ in the bottom half of both maximisers, so $c\in L_{\mathrm b}$,
and \eqref{eq:halvesgen} gives $\tau-c\in H_{\mathrm a}\subseteq\mathcal{S}$, whence
$\tau-c\le\max\mathcal{S}=m$, that is $\tau\le m+c$. The two displayed inequalities contradict each
other.

\emph{The minimum} is the mirror image. Put $m^{-}=\min\mathcal{S}$, let $i^{-}$ be its class,
$n=n_{i^{-}}\ge2$ and $c'=c_{i^{-},n-1}$, so that $m^{-}=c_{i^{-},n}$. Then
$m^{-}\in g_{\mathrm a}\cap g_{\mathrm b}$ if and only if $k_{i^{-}}=n-1$ in both, that is if and only
if both maximisers select \emph{every} increment of class $i^{-}$; again class $i^{-}$ cannot be tied,
since the maximiser that does not take the tied increment of a tied class omits one there. Hence every
increment of class $i^{-}$ exceeds $\tau$, in particular $\Delta_{i^{-}}(n-1)=m^{-}+c'>\tau$. And
$k_{i^{-}}=n-1$ places $c'$ in the top half of both, so $c'\in H_{\mathrm b}$, and
\eqref{eq:halvesgen} gives $\tau-c'\in L_{\mathrm a}\subseteq\mathcal{S}$, whence
$\tau-c'\ge\min\mathcal{S}=m^{-}$, that is $\tau\ge m^{-}+c'$. Contradiction.
\end{proof}

\end{document}
